\subsection{Resource States}

\begin{definition}
A \term{resource state} is a state, typically containing entanglement, on which a measurement pattern can be applied to conduct measurement-based quantum computation. For the purposes of this thesis, a resource state is a graph state.
\end{definition}

\begin{definition}
A \term{family} of resource states is some collection of resource states. A family of resource states $\Psi$ is \term{universal for MBQC} if, for every $n$-qubit state $\ket\phi$, there exists a resource state $\ket{\Psi} \in \Psi$ on at least $n$ qubits and a measurement pattern $\cB$ whose output is $\ket\phi$. We say a family of resource states is \term{$n$-qubit universal for MBQC} if this holds for all $m$-qubit $\ket\phi$ with $m \leq n$. When the meaning is clear, we may drop the phrase ``for MBQC'' and simply say a family of resource states is universal.
\end{definition}

We start by recalling some well-known results related to graph states.

\begin{fact}
\label{fact:graph-state-unique-construction}
Every graph state has a canonical construction by preparing $\ket+^{\otimes |V(G)|}$
and applying one $\CZ$ gate for each edge of $G$, uniquely up to reordering of the
commuting $\CZ$ gates.
\end{fact}

\begin{fact}
\label{fact:graph-circuit-depth-is-max-degree}
Let $G$ be a graph with maximum degree $\Delta$. Ignoring the initial preparation of $\ket{+}^{\otimes |V(G)|}$, $\ket{G}$ can be constructed using at most $\Delta+1$ layers of $\CZ$ gates.
\begin{proof}
Vizing's theorem states that any graph with maximum degree $\Delta$ admits an edge coloring with $\Delta + 1$ colors such that no two edges that share a vertex have the same color. All edges of the same color can be implemented in the graph state in the same layer of CZs. Thus, at most $\Delta + 1$ layers are needed.
\end{proof}
\end{fact}

\begin{fact}
\label{fact:vertex-deletion}
Given a graph state $\ket{G}$, a $Z$-basis measurement of a vertex is equivalent
to deleting that vertex, up to local Pauli byproducts on its neighbors.
\end{fact}

\begin{proof}
Consider a graph state $\ket{G}$ in its ZX-Calculus representation.
\begin{figure}[H]
\centering
\[
\begin{aligned}
%==================== LEFT DIAGRAM ====================
\begin{ZX}
    & & & & & & \\
    &
    \zxZ{} \ar[r] \ar[d] \ar[ul]
    & \zxH{} \ar[r]
    & \zxZ{} \ar[r] \ar[d] \ar[u]
    & \zxH{} \ar[r]
    & \zxZ{} \ar[d] \ar[ur]
    & \\
    &
    \zxH{} \ar[d]
    &
    & \zxH{} \ar[d]
    & \zxX{\alpha\pi}
    & \zxH{} \ar[d]
    & \\
    &
    \zxZ{} \ar[r] \ar[d] \ar[l]
    & \zxH{} \ar[r]
    & \zxZ{} \ar[r] \ar[d] \ar[ur]
    & \zxH{} \ar[r]
    & \zxZ{} \ar[d] \ar[r]
    & \\
    &
    \zxH{} \ar[d]
    &
    & \zxH{} \ar[d]
    &
    & \zxH{} \ar[d]
    & \\
    &
    \zxZ{} \ar[r] \ar[dl]
    & \zxH{} \ar[r]
    & \zxZ{} \ar[r] \ar[d]
    & \zxH{} \ar[r]
    & \zxZ{} \ar[dr]
    & \\
    & & & & & & 
\end{ZX}
&=
%==================== MIDDLE DIAGRAM ====================
\begin{ZX}
    & & & & & & \\
    &
        \zxZ{} \ar[r] \ar[d] \ar[ul]
        & \zxH{} \ar[r]
        & \zxZ{} \ar[r] \ar[d] \ar[u]
        & \zxH{} \ar[r]
        & \zxZ{} \ar[d] \ar[ur]
    & \\
    &
        \zxH{} \ar[d]
        &
        & \zxZ{\alpha\pi}
        &
        & \zxH{} \ar[d]
    & \\
    &
        \zxZ{} \ar[r] \ar[d] \ar[l]
        & \zxZ{\alpha\pi}
        &
        & \zxZ{\alpha\pi} \ar[r]
        & \zxZ{} \ar[d] \ar[r]
    & \\
    &
        \zxH{} \ar[d]
        &
        & \zxZ{\alpha\pi} \ar[d]
        &
        & \zxH{} \ar[d]
    & \\
    &
        \zxZ{} \ar[r] \ar[dl]
        & \zxH{} \ar[r]
        & \zxZ{} \ar[r] \ar[d]
        & \zxH{} \ar[r]
        & \zxZ{} \ar[dr]
    & \\
    & & & & & &
\end{ZX}
&=
%==================== RIGHT DIAGRAM ====================
\begin{ZX}
    & & & & & & \\
    &
    \zxZ{} \ar[r] \ar[d] \ar[ul]
    & \zxH{} \ar[r]
    & \zxZ{\alpha\pi} \ar[r]
    & \zxH{} \ar[r]
    & \zxZ{} \ar[d] \ar[ur]
    & \\
    &
    \zxH{} \ar[d]
    &
    &
    &
    & \zxH{} \ar[d]
    & \\
    &
    \zxZ{\alpha\pi} \ar[d]
    &
    &
    &
    &
    \zxZ{\alpha\pi} \ar[d]
    & \\
    &
    \zxH{} \ar[d]
    &
    &
    &
    & \zxH{} \ar[d]
    & \\
    &
    \zxZ{} \ar[r] \ar[dl]
    & \zxH{} \ar[r]
    & \zxZ{\alpha\pi} \ar[r]
    & \zxH{} \ar[r]
    & \zxZ{} \ar[dr]
    & \\
    & & & & & &
\end{ZX}
\end{aligned}
\]
\caption{Z plane measurements correspond to vertex deletion.}
\label{fig:grid-zx-rewrite}
\end{figure}
By Lemma \ref{lemma:measurement-plane-maps}, the ZX effect of a Z plane measurement is a $\zx{\zxX{a \pi}}$ where $a$ is the measurement outcome bit. By \zxref{cp}, we may push the $\zx{\zxX{a \pi}}$ through the vertex being measured, deleting it in the process. The result is a graph state with the measured vertex removed, and a local $Z^{a}$ byproduct on each neighbor of the measured vertex.
\end{proof}

\begin{fact}
\label{fact:can-local-complement-graph-state}
Given a graph state $\ket{G}$, applying $\sqrt{X} = \zx{\rar & \zxX{\pi/2} \rar &}$ on a vertex $v$ and $\sqrt{Z} = \zx{\rar & \zxZ{\pi/2} \rar &}$ up to global phase on $N_G(v)$ is equivalent to a local complementation around $v$, and yields $\ket{G \star v}$. Applying a $Y$ plane measurement on $v$ instead gives $\ket{G \star v - v}$ with additional phases.
\end{fact}

\begin{proof}
The first claim is Lemma \ref{lemma:zx-local-complement-without-deletion} and the second is Lemma \ref{lemma:remove-half-pi-from-graph}.
\end{proof}

\begin{fact}
\label{fact:can-pivot-graph-state}
Given a graph state $\ket{G}$ and an edge $uv \in E(G)$, applying $\Had$ on vertices $u$, $v$ and $Z$ on $N_G(u) \cap N_G(v)$ gives $\ket{G \land uv}$. Applying a $X$ plane measurement on both $u$ and $v$ instead gives $\ket{G \land uv - u - v}$ with additional phases. 
\end{fact}

\begin{proof}
The first claim is Lemma \ref{lemma:zx-pivot-without-deleting} and the second is Lemma \ref{lemma:remove-pi-from-graph}.
\end{proof}

\begin{definition}
\label{def:grid}
A $m \times n$ grid graph $G = (V,E)$ has
\begin{align}
    V &= \{(i,j) \mid i \in [m], j \in [n]\}, \\
    E &= \{\{(i_1,j_1),(i_2,j_2)\} \mid (i_1,j_1),(i_2,j_2) \in V, |i_2-i_1| + |j_2-j_1| = 1\}.
\end{align}
In other words, $m$ rows and $n$ columns with edges between nodes of taxicab distance $1$. Denote such graphs as $\Grid_{m \times n}$. Let \term{$\Grid$} denote the family of all such graphs.
\end{definition}

\begin{theorem}
\label{thm:grid-universal}
The $\Grid$ is universal.
\end{theorem}

\begin{figure}[!htbp]
\centering
\[
\begin{ZX}
% row 1: top padding
& & & & & & & & & & & & \\
% row 2: y = 1
&
\zxZ{} \ar[d] \ar[u]
& & & & & &
\zxZ{} \ar[d] \ar[u]
& & & &
\zxZ{} \ar[d] \ar[u]
& \\
% row 3: between y = 1 and y = 2
&
\zxH{} \ar[d]
& & & & & &
\zxH{} \ar[d]
& & & &
\zxH{} \ar[d]
& \\
% row 4: y = 2
&
\zxZ{} \ar[r] \ar[d] \ar[l]
& \zxH{} \ar[r]
& \zxZ{} \ar[r] \ar[u]
& \zxH{} \ar[r]
& \zxZ{} \ar[r] \ar[u]
& \zxH{} \ar[r]
& \zxZ{} \ar[d] \ar[r]
& & & &
\zxZ{} \ar[d] \ar[r]
& \\
% row 5: between y = 2 and y = 3
&
\zxH{} \ar[d]
& & & & & &
\zxH{} \ar[d]
& & & &
\zxH{} \ar[d]
& \\
% row 6: y = 3
&
\zxZ{} \ar[d] \ar[l]
& & & & & &
\zxZ{} \ar[d] \ar[r]
& & & &
\zxZ{} \ar[d] \ar[r]
& \\
% row 7: between y = 3 and y = 4
&
\zxH{} \ar[d]
& & & & & &
\zxH{} \ar[d]
& & & &
\zxH{} \ar[d]
& \\
% row 8: y = 4
&
\zxZ{} \ar[d] \ar[l]
& & & &
\zxZ{} \ar[r] \ar[d] \ar[l]
& \zxH{} \ar[r]
& \zxZ{} \ar[r]
& & & &
\zxZ{} \ar[d] \ar[r]
& \\
% row 9: between y = 4 and y = 5
&
\zxH{} \ar[d]
& & & &
\zxH{} \ar[d]
& & & & & &
\zxH{} \ar[d]
& \\
% row 10: y = 5
&
\zxZ{} \ar[d] \ar[l]
& & & &
\zxZ{} \ar[d] \ar[l]
& & & & & &
\zxZ{} \ar[d] \ar[r]
& \\
% row 11: between y = 5 and y = 6
&
\zxH{} \ar[d]
& & & &
\zxH{} \ar[d]
& & & & & &
\zxH{} \ar[d]
& \\
% row 12: y = 6
&
\zxZ{} \ar[d] \ar[l]
& & & &
\zxZ{} \ar[r] \ar[d] \ar[l]
& \zxH{} \ar[r] 
& \zxZ{} \ar[r] \ar[d]
& \zxH{} \ar[r]
& \zxZ{} \ar[r] \ar[d]
& \zxH{} \ar[r]
& \zxZ{} \ar[d] \ar[r]
& \\
% row 13: between y = 6 and y = 7
&
\zxH{} \ar[d]
& & & &
\zxH{} \ar[d]
& & & & & &
\zxH{} \ar[d]
& \\
% row 14: y = 7
&
\zxZ{} \ar[d]
& & & &
\zxZ{} \ar[d]
& & & & & &
\zxZ{} \ar[d]
& \\
% row 15: bottom padding
& & & & & & & & & & & & 
\end{ZX}
\]
\caption{A $\Grid$ graph state with some vertices removed, creating three vertical ``wires'' for a three qubit computation in the MBQC model. The horizontal connections between wires may entangle the qubits, given the right measurements.}
\label{fig:padded-zx-diagram}
\end{figure}

\begin{proof}
The full proof is quite involved, so we only summarize the main ingredients. The original proof is given in \cite{MBQCOriginal}, and a more accessible proof is available in \cite{Raussendorf_2003}. A modern proof through the ZX-Calculus is available in Chapter 9 of \cite{PQS}. 

There are three great difficulties in the proof. We informally discuss them. First, we need to be able to implement arbitrary single-qubit gates. This is done by deleting (via $Z$ plane measurements) vertices to create line graphs, each resulting in a chain like those in Example \ref{ex:teleportation-gadget}. It can be shown that this allows the application of arbitrary single-qubit gates up to some Pauli error.

Second, we need to be able to entangle qubits. If the two line graphs containing the qubits we wish to entangle are adjacent, i.e. there are no other line graphs between them, then an orthogonal connecting chain of vertices can be used to implement an
entangling operation between the two logical wires. Now, if they are not adjacent, we can swap qubits between lines (by implementing three CNOT gates) until they become adjacent.

Third, we need to ensure that measurements are performed in the correct order. For example, suppose we delete all vertices connecting the input and output vertices. Then this computation is deterministic and not invertible, and thus does not implement a unitary. As it turns out, there is a criterion called \emph{flow} that ensures this does not happen, as well as variants of flow such as \emph{gflow} and \emph{Pauli flow} that further ensure that we can deal with the Pauli errors. Interested readers should see \cite{McElvanney_2023,Backens_2021,PQS} for more details.
\end{proof}

